供应链映射的核心不是把所有公司都排进一张长表，而是回答三件事：公司处在上游、中游还是下游；它的产品是不是 AIDC 的真实瓶颈；这个瓶颈能不能转成 2026E 收入、毛利和 EPS。下游 AI 资本开支只能证明行业方向，不能自动证明单一供应商收入增长。本版已经把 59 个候选逐一穿透到七类字段：收入暴露、客户/平台、订单/交付、产能/认证、ASP/价格代理、利用率/良率代理和毛利影响；直接证据进入基础分母，代理或模型代理只限制模型边界，需求锚只能验证方向，OCR 噪声和来源耗尽字段不入模。完整 58 个核心候选公司级矩阵移至附录，本章正文只解释价值链怎样传导、哪些环节能转成 EPS、哪些只能保留观察。

\textbf{算力与存储：价值量最大，但 A 股映射必须折价。} GPU/AI ASIC、CPU、HBM/DRAM、DPU/NIC、BMC、交换 ASIC 和高速 SerDes 决定集群吞吐、显存带宽、通信延迟和整柜功耗，是 AIDC 成本与性能的最上游约束。但这个环节的利润池并不会自然流向 A 股，原因是核心 GPU、HBM、先进封装和高端交换芯片仍主要由海外供应链主导，A 股公司更多暴露在国产算力芯片、CPU、存储模组、接口芯片、EDA/IP、测试设备和服务器零部件的间接位置。因此，上游技术越核心，不等于估值信用越高；只有当公司能拿出产品销售、平台适配、客户交付、毛利率和现金转化证据时，才可进入盈利模型。否则，NVIDIA、AMD、博通、SK 海力士、美光等只能作为需求锚或技术路线锚，不能替代 A 股公司的收入确认。

\textbf{服务器、整柜与网络设备：收入弹性最大，利润率约束最硬。} 算力芯片、HBM、光模块、PCB、连接器、电源、散热和高速交换芯片最终需要被集成为 AI 服务器、整柜系统、交换机和集群网络。这个环节把上游 BOM 变成整机收入，也是 A 股最容易看到收入弹性的地方，工业富联、浪潮信息、中科曙光、紫光股份、锐捷网络等都处在这条线上。问题在于整机和 ODM/OEM 业务天然毛利率偏薄，且受客户议价、核心芯片配额、库存、应收账款和交付节奏影响。研究时不能只看 AI 服务器收入增速，还要看产品结构、在手订单、客户集中度、存货周转、经营现金流和毛利率是否同步改善；否则收入增长可能只是在替上游芯片和客户账期承担资本占用。

\textbf{光通信：最清晰利润池，也是最容易被提前定价的环节。} AI 集群从单机训练走向万卡、十万卡网络后，800G/1.6T 光模块、DSP、EML、硅光、FAU、薄膜铌酸锂、CPO/LPO/NPO 等技术路线决定东西向流量的带宽、功耗和成本。A 股及港股可跟踪公司集中在光模块、光器件和上游材料设备，盈利兑现能力明显强于多数概念环节。中际旭创、新易盛、天孚通信、光迅科技、源杰科技、联特科技、太辰光等不能放在同一个估值框里：龙头模块厂要看海外客户分配、800G/1.6T ASP、良率和产能爬坡；器件和芯片公司要看是否进入龙头供应链、单价占比和替代节奏；CPO/LPO 相关标的则要先证明从技术路线到订单的闭环。这里可以给较高盈利权重，但必须用客户、价格、良率和订单做季度复核。

\textbf{PCB、CCL 与电子材料：从普通电子周期切换为高端结构占比逻辑。} AI 服务器、交换机、UBB、OAM、GPU baseboard 和高速背板对高多层、高速低损耗材料、HVLP 铜箔、阻抗控制、散热和良率提出更高要求，PCB/CCL 不再只是传统电子景气 beta。沪电股份、胜宏科技、深南电路、生益科技、华正新材、金安国纪等公司的关键不是“有没有 AI 概念”，而是高端产品收入占比、产能认证、客户导入、良率、材料成本传导和单机价值量。PCB/CCL 的利润池可验证性较强，因为它最终会落在收入结构、毛利率和产能利用率上；但如果扩产先行而客户认证、良率或高端订单没有同步兑现，估值必须按周期股而不是成长股处理。

\textbf{供配电与液冷：从配套件变成 AIDC 交付门槛。} 高功率 AI 机柜把传统 IDC 的电力和散热瓶颈前置。UPS、HVDC、变压器、PDU、母线、CDU、冷板、Manifold、泵阀、换热器、氟泵和运维监控不再只是机房配套，而是决定能否上架、能否稳定运行、能否通过客户验收的交付门槛。英维克、科华数据、申菱环境、佳力图、依米康、同飞股份、飞荣达、科士达、麦格米特、盛弘股份等公司应按“认证--订单--交付--验收--毛利率--回款”逐级判断，而不是按液冷或电源标签一刀切。这个环节的风险是项目制收入确认和客户验收节奏，毛利率也会受定制化、原材料、工程交付和售后维护影响；因此只有批量交付和现金流改善同时出现，才应给目标价上修。

\textbf{AIDC/IDC 运营：不是设备 beta，而是资产周转和利用率模型。} 下游运营商把土地、电力、网络、机柜、冷却和客户租约转成现金流，核心变量是 MW、机柜数、上架率、PUE、电价、折旧、融资成本、客户合同期限和回款质量。润泽科技、奥飞数据、光环新网、宝信软件、数据港、杭钢股份等公司的研究框架与设备公司完全不同：新增项目和算力园区只能说明资产规模，真正决定 EPS 的是客户上架速度、单位电力成本、机柜租金、折旧摊销和资本开支回收期。AIDC 运营商可以享受 AI 需求重估，但如果利用率或现金回款不能跟上，估值应被折回 REIT/公用事业式现金流模型，而不是设备成长股模型。

由此得到的投资结论是：AIDC 产业链不是越上游越值得买，也不是表格里字段越多就越有投资价值。上游芯片和存储决定真实技术瓶颈，但 A 股盈利穿透弱；服务器和网络设备最容易放大收入，但利润率、库存和应收账款是硬约束；光模块和 AI PCB/CCL 是当前最清晰的利润池，但需要持续验证 ASP、良率、客户分配和产品结构；液冷、供配电和 AIDC 运营提供二阶增量，但必须跟踪认证、订单、验收、上架率和现金流。正文的估值分层正是从这条因果链出发，而不是从概念归类出发。

\begin{exhibitbox}[成本拆解、利润池与 26E 验证变量]
\scriptsize
\begin{longtable}{L{1.8cm}L{3.2cm}L{3.7cm}L{3.2cm}L{3.6cm}}
\toprule
\textbf{位置} & \textbf{成本/价值量构成} & \textbf{利润池读法} & \textbf{26E 验证变量} & \textbf{A股映射含义}\\
\midrule
\endhead
算力与存储 & GPU/AI ASIC、CPU、HBM/DRAM、SSD、DPU/NIC、交换 ASIC & 海外芯片和 HBM 占主导；A 股多为国产替代、接口/存储/平台映射 & 芯片 ASP、HBM 供给、先进封装产能和国产平台验证 & 多数 A 股节点低纯度；客户/收入直接口径不足时只给期权或观察信用\\
服务器/整柜 & GPU/CPU/HBM、PCB、电源、连接器、铜缆、结构件、液冷部件 & ODM 收入规模大但毛利率薄，利润来自交付效率、BOM 管理和客户份额 & AI 服务器/交换机订单、毛利率、存货和应收账款周转 & 工业富联、浪潮信息、中科曙光、紫光股份等需要拆 AI 服务器纯度\\
网络与光互联 & DSP、EML/硅光、FAU、透镜、陶瓷、PCB、测试设备 & 800G/1.6T 模块和精密器件是高弹性利润池，ASP 与良率决定 EPS & 客户分配、产品代际、ASP、毛利率、订单持续性 & 中际旭创/新易盛盈利兑现更强，光器件/CPO 链需折价看待\\
PCB/材料 & 高速 CCL、铜箔、树脂、玻纤布、钻孔/电镀设备 & 价值从普通 PCB 转向高多层、HDI、UBB、高速交换机板和低损耗材料 & 高端产品结构、扩产爬坡、良率、材料成本传导 & 沪电/胜宏/深南/生益的估值核心是高端结构占比，不是普通 PCB 周期\\
供配电 & 变压器、UPS/HVDC、PDU、母线、BBU、电池、开关柜 & AIDC 功率密度提升把电力系统从配套件变成交付约束 & 认证、项目中标、交付验收、毛利率和经营现金流 & 科华、金盘等必须证明数据中心项目收入而非泛电力设备叙事\\
液冷/温控 & CDU、冷板、Manifold、泵阀、冷却液、精密空调、冷源 & 液冷价值取决于设计导入、系统可靠性和售后服务，不是单一部件概念 & 客户认证、批量交付、验收、液冷收入占比、毛利率 & 英维克/申菱/高澜等要用项目和收入拆分验证估值信用\\
AIDC/IDC 运营 & 土地/能耗指标、电网接入、服务器网络设备、冷却系统、融资 & 资产端看 MW 和机柜，利润端看上架率、电价、折旧、客户租约 & 新增 MW、上架率、单位电力成本、租约久期、经营现金流 & 润泽/奥飞/光环等不是设备 beta，核心是资产周转和利用率\\
\bottomrule
\end{longtable}
\sourcenote{analysis/value\_chain\_economics.md；analysis/chain\_business\_research.md。26E 验证变量是 AStock 模型所需字段，不是外部一致预期。}
\end{exhibitbox}

上表只用于压缩呈现成本、利润池和验证变量，不能替代正文判断。真正需要被解释的是每个环节的“价值量--收入确认--利润率--现金流--估值信用”链条：上游需求锚只能说明方向，中游瓶颈要落到出货和毛利，下游资产要落到上架率和现金回收。完整 58 个核心候选的上游业务、下游业务、业务关联、核心技术、核心营收业务和 2026E 预期已经作为附录证据索引披露；正文不再逐行铺表，而是按产业链因果关系给出可投资性判断。

上游位置的关键矛盾是“技术核心”和“A 股可证伪收入”并不总一致。GPU、HBM、交换 ASIC、DSP、CPO 核心 ASIC、ABF 等是真正高价值环节，但 A 股直接映射弱，更多是国产替代或设备/材料/平台的间接暴露。中游位置更适合做盈利预测：服务器/交换机 ODM 看收入规模和毛利率，光模块看 800G/1.6T 出货、ASP 和良率，PCB/CCL 看高端产品结构和材料成本传导，液冷/供配电看认证转订单、项目交付和验收。下游位置则更像资产运营模型：AIDC 运营商看 MW、上架率、电价、折旧和租约，而不是单纯看 AI 主题热度。

核心技术也必须分层看。光模块的技术焦点在高速光电设计、封装、热管理、耦合精度和测试良率；PCB 的技术焦点在低损耗材料、高多层叠构、阻抗控制和良率；液冷的技术焦点在冷板/CDU/Manifold 系统可靠性、漏液防护和维护；供配电的技术焦点在功率转换效率、冗余可靠性和高密模块交付；AIDC 运营的技术焦点则是机房电力、网络、冷却、SLA 和调度。用同一个 PE 去套这些业务没有意义，必须先判断收入驱动和利润池在哪里。

\begin{exhibitbox}[核心标的上下游位置、核心技术与收入口径]
\scriptsize
\begin{tabularx}{\textwidth}{L{2.0cm}L{2.2cm}X X X}
\toprule
\textbf{环节} & \textbf{覆盖标的} & \textbf{上游输入} & \textbf{核心技术/营收} & \textbf{下游位置与26E验证}\\
\midrule
光模块/光器件 & 中际旭创、新易盛、天孚通信 & DSP、EML/VCSEL、硅光、FAU、透镜、陶瓷、PCB、测试设备 & 800G/1.6T 光模块、光引擎、精密无源器件；核心技术是高速封装、热管理、耦合精度和良率 & 下游为 AI 训练/推理集群和交换机端口；26E 验证 ASP、客户分配、良率、毛利率和订单持续性\\
AI PCB/CCL & 沪电股份、胜宏科技、深南电路、生益科技 & 高速 CCL、铜箔、树脂、玻纤布、钻孔/电镀设备 & 高多层板、HDI、UBB、交换机/路由器板和低损耗材料；核心技术是叠构、阻抗控制和材料匹配 & 下游为 AI 服务器、HPC 和交换机；26E 验证高端产品结构、扩产爬坡、良率和材料成本传导\\
AI 服务器/网络设备 & 工业富联、浪潮信息、中科曙光、紫光股份 & GPU/AI ASIC、CPU、HBM、PCB、电源、结构件、光模块 & 服务器、整柜系统、高速交换机、国产算力平台；核心技术是系统集成、BOM 管理和供应链交付 & 下游为海外 CSP、云厂商、运营商和政企智算；26E 验证订单转收入、毛利率、库存和应收\\
液冷/温控 & 英维克、申菱环境 & 压缩机、泵阀、冷板、CDU、Manifold、冷却液、精密空调 & CDU、冷板、液冷机柜、冷源和精密空调；核心技术是可靠性、漏液防护、换热效率和维护 & 下游为服务器厂、IDC、云厂和运营商；26E 验证客户认证、批量交付、验收、收入占比和毛利率\\
供配电/能源 & 科华数据、金盘科技 & 功率器件、变压器、电池、母线、开关柜、硅钢、铜材 & UPS/HVDC、预制电力模组、干式变压器和数据中心供电设备；核心技术是效率、冗余和认证 & 下游为 AIDC/IDC、运营商和园区；26E 验证认证转订单、项目交付、验收、毛利率和现金流\\
AIDC/IDC 运营 & 润泽科技、奥飞数据、光环新网 & 土地/能耗指标、电网接入、服务器网络设备、冷却系统、融资 & AIDC/IDC 托管、算力服务、边缘算力和云服务；核心技术是 MW 交付、上架率、SLA 和调度 & 下游为云厂、模型公司、互联网和政企客户；26E 验证新增 MW、上架率、电力成本、折旧和租约\\
\bottomrule
\end{tabularx}
\sourcenote{data/supply\_chain\_relationships.json；data/growth\_driver\_model.json；data/current\_valuation\_model\_20260630.json；data/core\_candidate\_extended\_valuation\_model\_20260701.json。逐公司 56 个目标价/公允价值模型与 58 个核心候选公司卡片保存在 analysis/core\_candidate\_company\_cards.md。}
\end{exhibitbox}

本报告给予较高证据权重的关系包括：工业富联 AI 服务器与高速交换机收入披露、沪电股份数据通信 PCB 与 AI 服务器/HPC 分项披露、英维克液冷产品线和客户案例、科华数据高密 UPS 认证、润泽科技 AIDC 项目和业务展望。它们不是因为“属于 AIDC 概念”而进入核心池，而是因为至少有一条从产品到收入或项目的证据链。相反，HBM、BMC、DPU/NIC、DSP、CPO/NPO 核心 ASIC、ABF、高端 HVLP 铜箔、绿电和通用 EPC 虽然属于真实链条，但多数 A 股映射缺高纯度或缺直接收入证据，只能做观察或条件估值。

增长预期上，2026E 不能简单写成“AI 高景气”。服务器和交换机要看订单能否转收入且毛利率不塌；光模块要看 800G/1.6T 的 ASP、客户分配和良率；PCB/CCL 要看高端产品结构能否覆盖扩产和材料成本；液冷/供配电要看认证是否进入批量交付；AIDC 运营要看新增 MW、上架率、电力成本和折旧压力。后续季度若这些变量没有兑现，估值模型必须下修，而不是用远期 TAM 继续支撑目标价。
